\documentclass{article}
\usepackage{fontspec}
\setmainfont{Times New Roman}
\usepackage{amsmath}
\usepackage{amssymb}
\usepackage{amsfonts}

\begin{document}

\section{Допълнение на събитие}

Много често пространството от елементарни събития $\Omega$ си го представяме като множество, независимо дали е дискретно или не. Нека имаме събитие $A$, което е подмножество на $\Omega$, тоест:
\[
A \subseteq \Omega
\]

\subsection{Дефиниция}
Допълнението на събитието $A$, което бележим с $A^c$, се дефинира като разликата на пространството $\Omega$ и множеството $A$:
\[
A^c = \Omega \setminus A
\]
С други думи, ако едно елементарно събитие $\omega$ принадлежи на допълнението $A^c$, то не принадлежи на $A$:
\[
\text{ако } \omega \in A^c, \text{ то } \omega \notin A
\]
Самото пространство $\Omega$ се състои от събитието $A$, обединено с неговото допълнение:
\[
\Omega = A \cup A^c
\]

\subsection{Примери за допълнение}
Ако имаме някакъв признак, по който определяме множеството $A$, то допълнението $A^c$ съдържа всички елементарни събития, които не притежават този признак.
\begin{itemize}
    \item \textbf{Пример 1 (Пол):} Ако събитието $A$ е множеството на жените, то допълнението на $A$ ще бъде множеството на мъжете:
    \[
    A^c = \{\text{мъже}\}
    \]
    \item \textbf{Пример 2 (Хвърляне на зарче):} Ако разгледаме хвърляне на зарче и събитието да се падне нечетно число, то неговото допълнение са четните числа:
    \[
    A^c = \{2, 4, 6\}
    \]
    \item \textbf{Пример 3 (Гласуване):} Ако $A$ е събитието „гласували за партия Х“, то допълнението $A^c$ е множеството на „негласувалите за партия Х“.
\end{itemize}

\section{Свойства на операциите със събития}

Операциите с множества притежават някои естествени свойства.
\begin{itemize}
    \item \textbf{Комутативност:} Обединението и сечението на две събития не зависят от реда, в който ги записваме:
    \[
    A \cup B = B \cup A
    \]
    \[
    A \cap B = B \cap A
    \]
    \item \textbf{Асоциативност:} Ако имаме три събития $A$, $B$ и $C$, можем да ги групираме по различен начин при обединение:
    \[
    (A \cup B) \cup C = A \cup (B \cup C)
    \]
\end{itemize}

\end{document}